\documentclass[12pt]{article}
\usepackage[utf8]{inputenc}
\usepackage[T1]{fontenc}
\usepackage{hyperref}
\title{Ditammari Language Resources for K-12 Curriculum Development}
\author{Compiled by Claude Sonnet 5 (claude-sonnet-5) for Andr\'as Kornai}
\date{}
\begin{document}
\maketitle

\section{Introduction}
Ditammari (Tammari, Otammari) is spoken in Benin and Togo. Sources disagree on the number of speakers; about 150,000 is the figure given for this survey. The search found no dedicated primer, but it did show that the language has a Latin orthography and a printed Bible.

\section{Primary Educational Resources}
No specific textbook or primer was found in search. Omniglot says Ditammari is taught in some primary schools and used in some radio programmes: \url{https://www.omniglot.com/writing/tammari.htm}. An Italian prosody paper says the written variety has an orthography and school texts: \url{https://www.aisv.it/PubblicazioniAISV/II_AISV/presentations/Prosodia/Romano%20A.-Mouti%20O.pdf}

\section{Linguistic and Reference Materials}
\begin{itemize}
\item Wikipedia on the Tammari language, with the alphabet (tone marked by accents, nasalization by tildes): \url{https://en.wikipedia.org/wiki/Tammari_language}
\item Selective bibliography on Beninese languages and literacy (CNLA, Cotonou, 1990), indexed by IDRC. I did not check whether it lists Ditammari titles: \url{https://idl-bnc-idrc.dspacedirect.org/handle/10625/36911?show=full}
\end{itemize}

\section{Religious and Community Resources}
A Ditammari Bible (\emph{Kuyie Nn\'a\'a\v{n}t\`i P\'at\'iri}), published in Cotonou by the Alliance biblique du B\'enin, is cited in the Wikipedia article above. I did not find a direct URL for it.

\section{Digital and Audio Resources}
Omniglot mentions radio programmes. No audio archive was found.

\section{Curriculum Development Context}
There is an established orthography and some school use, but no documented curriculum materials turned up. The next steps would be to contact the Alliance biblique du B\'enin and the national literacy bodies in Benin.
\end{document}